\chapter{系统设计}
\label{ch:design}

\section{总体目标}
系统需要支持初始化、插入、查询、删除、批量加载、结构遍历和指标采集；同时要保留理论结构中的关键接口，使实验能够观察到探针、踢出、元数据
位数、tier 分布和扩缩容事件。

从复杂度目标看，理想模型希望查询在局部路由与元数据访问均为常数时间时达到 \(O(1)\)，插入和删除以 \(O(k)\) 次交换为主要代价，总额外空间接近 \(O(n\log^{(k)}n)\) 比特\cite{pagh2021otsh}。当前系统更强调可解释和可验证：它保留 Facility、Cubby、Router、MiniArray、商压缩、tail 与后台重建等结构，但在若干微观部件上采用工程化实现。因此本文后续讨论会区分“理论目标”和“当前实现”。

系统采用 \texttt{Facility -> Cubby -> Slot} 的三层结构。这个结构方便把宏观容量管理、局部插入策略和每槽元数据分开处理。

\subsection{Facility：设施级路由与分布管理}
Facility 是全表的宏观分片。设全局容量参数为 \(N\)，单个设施管理约 \(K=\mathrm{polylog}\,N\) 规模的局部元素。对键 \(x\) 先计算可逆置换 \(y=\pi(x)\)，再从 \(g^\ast(x)=y\bmod N\) 中拆出设施编号 \(r\) 和设施内桶编号 \(b\)。系统访问第 \(r\) 个 Facility 的 \texttt{D[b]}，由该局部路由器返回目标 Cubby 和槽位。

一个 Facility 内部包含三类状态：
\begin{itemize}
  \item \texttt{D}：长度为 \(K\) 的路由器数组，每个元素管理一个小规模键集合到 \((\mathrm{Cubby},\mathrm{slot})\) 的映射；
  \item 多 tier cubby 集合：不同 tier 表示不同容量尺度，用于近似理论中的分层 cubby；
  \item tail：唯一允许非满的 cubb。
\end{itemize}

tail 的存在使“除 tail 外尽量保持满”的不变量更容易维护。插入优先进入 tail；删除非 tail 元素后，可从 tail 搬迁一个元素补洞。这样做可以
减少普通 cubby 长期半空造成的空间浪费，也为后续重建提供较稳定的输入，为装填因子可达到 1 提供了可能。

\subsection{Cubby：微观存储单元}
Cubby 是实际保存元素的基本单元。逻辑上，一个 Cubby 对应理论中的局部 k-kick 实例；工程上，它包含以下状态：
\begin{itemize}
  \item \texttt{slots}：槽数组，保存商压缩后的 payload，而不是简单保存完整键；
  \item \texttt{occupied}：记录已占用槽位，便于迁移、删除和结构导出；
  \item \texttt{FreeSlotTree}：用于快速寻找空槽；
  \item \texttt{MiniArray meta}：按槽保存变长 \texttt{MetaEntry} 编码。
\end{itemize}

在理论构造中，Cubby 内部应按 k-kick 树的分层优先级执行插入。当前系统保留踢出链和移动计数，用有限随机探针和局部搬迁近似该行为。这样可以
让插入、删除、查询路径完整运行，并能通过 \texttt{kick\_count} 与移动次数观察更新成本，但它不是论文中 packed 微观结构的逐行复现。

\subsection{Slot 与 MetaEntry：从落槽到键恢复}
每个槽位不仅保存 payload，还需要保存恢复原键所需的信息。系统把这些信息压入 \texttt{MetaEntry}，主要包括：
\begin{itemize}
  \item \texttt{delta}：偏好槽 \(g_I(x)\) 与实际槽 \(i\) 的差值；
  \item \texttt{mid\_bits}：从 Cubby 内位置拼回设施内路由低位所需的中间位；
  \item \texttt{router\_bits}：路由辅助信息；
  \item \texttt{insert\_bits}：插入过程相关标签。
\end{itemize}

\texttt{MetaEntry} 的意义在于把“元素为什么能在这个槽被找到”压缩成少量可编码字段。查询时，系统并非只根据键重新线性探测，而是先经 Router 定
位到槽，再利用 payload 与 MetaEntry 恢复通过反函数 \(\pi^{-1}(y)\) 得到的原键，最后用逆置换校验原键。

\section{紧凑构件设计}
\subsection{MiniArray}
MiniArray 用来保存变长元数据。理想结构可理解为单字宽节点组成的 packed B-树：叶子保存实际比特，内部节点保存偏移与枢轴，配合查找表实现均摊常数时间访问和更新。它解决的问题是：如果每个槽的元数据长度不同，普通定长数组会浪费空间，普通动态字节串又难以快速定位。

当前系统采用较直接的工程版本：用连续 bit buffer 保存所有条目，用 \texttt{offsets} 和 \texttt{bitlens} 定位每个槽的编码。更新某个槽时允许重新打包整个 Cubby 的元数据。由于 Cubby 规模被限制在 polylog 范围内，这个代价在原型系统中可接受；同时 \texttt{bits\_total} 能直接作为空间统计的输入，方便实验章估计元数据开销。

\subsection{Local Query Router}
Local Query Router 负责少量键的快速定位。设计上，它把键哈希成随机二进制串，构造压缩前缀 Trie，只保存能区分键集合的最短前缀，并在叶子处记录目标 Cubby 与 slot。Trie 的编码长度随局部元素数增长，理论上可以结合查表做到常数时间跳转。

当前系统的 \texttt{Router} 使用 64 位哈希串上的二叉 trie，并处理哈希碰撞桶。插入和删除之后允许重建 trie 与编码长度统计。这样的策略牺牲了理论中最细粒度的 \(O(1)\) 更新保证，但代码简单、状态清楚，且局部元素数较小时开销可控。实验上，\texttt{router\_probe\_steps} 和 \texttt{bits\_total} 可用来观察路由成本和编码规模。

\section{商压缩流程}
商压缩的目标是避免在槽中重复存储完整键。设 \(\pi\) 是 64 位可逆置换，\(g^\ast(x)=\pi(x)\bmod N\)。\(g^\ast(x)\) 的低位已经用于选择 Facility、路由桶和 Cubby 内偏好位置，因此这些位可以部分由寻址上下文隐含携带。Cubby 的 \texttt{slots} 保存的是去掉低 \(\log N\) 位后的 quotient payload；\texttt{MetaEntry} 保存把实际槽位还原为偏好槽、再拼回 \(g^\ast(x)\) 所需的信息。

查询流程因此分为两段。第一段是定位：由 \(\pi(x)\) 的低位进入对应 Facility 和 Router，得到目标 Cubby 与 slot。第二段是校验：从 slot 读取 payload 和 MetaEntry，根据设施编号、Cubby 容量、slot 索引和 \texttt{delta}/\texttt{mid\_bits} 重建 \(\pi(x)\)，再调用 \(\pi^{-1}\) 得到原键。若恢复出的键与查询键一致，则返回命中。这个设计把一部分键位从显式存储转移到寻址过程和元数据中，是系统接近信息论空间目标的主要来源。

\section{动态重建与扩缩容}
论文的构造要求不同 tier 的 cubby 数量维持在目标分布附近。设计文档中使用近似目标
\[
  t_j \approx
  \frac{\bigl(\log^{(j)} n\bigr)^2}
       {\bigl(\log^{(j+1)} n\bigr)^2}.
\]
当某一级 cubby 过少时，可以把更高一级 cubby 拆分；当某一级 cubby 过多时，可以把若干低级 cubby 合并。重建本身需要搬迁元素，如果一次性完成会造成延迟尖峰，因此系统把重建任务放入调度队列，每次更新只消耗固定预算。

全局扩缩容采用双表迁移。活动表负责新请求；当元素数超过负载上界或低于容量下界时，系统创建新 \texttt{TableState}，旧表保留为迁移来源，随后按设施逐步搬迁元素。这样做避免一次扩容搬动所有元素，也便于通过 \texttt{resize\_start} 和 \texttt{resize\_finish} 等指标观察迁移过程。

\section{主操作路径}
\textbf{插入路径。} 系统先计算 \(y = \pi(x)\)，拆出设施和桶编号 \(r\) 和 \(b\)，在 \texttt{D[b]} 中检查键是否已存在。若不存在，将元素
放入该 Facility 的 tail Cubby，选择槽位并写入 quotient payload 与 MetaEntry，然后更新空闲槽结构和 Router。如果 tail 满，则创建新
的 tail 或触发局部重建；全局规模变化后再检查是否需要扩容。

\textbf{查询路径。} 查询同样从 \(\pi(x)\) 出发定位 Facility 和 Router。Router 返回 Cubby 与 slot 后，系统读取 slot payload 和 
MetaEntry，反函数 \(x^{`} = \pi^{-1}(y)\) 。该路径的关键不在于遍历长探针序列，而在于局部路由和元数据能否把目标槽直接指出来。

\textbf{删除路径。} 删除先复用查询定位目标槽，再清空 slot、MetaEntry 和 Router 记录。如果删除发生在非 tail cubby，系统从 tail 迁移
一个元素填补空洞，保持普通 cubby 的满载倾向。删除后若规模过低，则触发缩容迁移。

\section{本章小结}
本章给出了当前系统的总体设计。Facility 负责宏观路由和分布管理，Cubby 负责局部存储和受控移动，Slot 与 MetaEntry 负责商压缩下的键恢复，MiniArray 和 Router 负责紧凑元数据与局部定位。下一章将结合代码文件说明这些设计如何具体实现，以及哪些部分属于工程化近似。